% Generated draft: requires mathematical and source review.
\documentclass[11pt]{article}
\usepackage[margin=25mm]{geometry}
\usepackage{fontspec}
\setmainfont{Latin Modern Roman}
\usepackage{amsmath,amssymb,mathtools,unicode-math}
\setmathfont{Latin Modern Math}
\usepackage{enumitem,needspace,xurl,hyperref}
\setlength{\parindent}{0pt}
\setlength{\parskip}{5pt}
\setlength{\emergencystretch}{3em}
\newcommand{\N}{\mathbb{N}}
\newcommand{\Z}{\mathbb{Z}}
\newcommand{\Q}{\mathbb{Q}}
\newcommand{\R}{\mathbb{R}}
\newcommand{\C}{\mathbb{C}}
\newcommand{\F}{\mathbb{F}}
\newcommand{\A}{\mathbb{A}}
\newcommand{\PP}{\mathbb{P}}
\newcommand{\E}{\mathbb{E}}
\newcommand{\eps}{\varepsilon}
\newcommand{\dd}{\,\mathrm{d}}
\begin{document}
\section*{UnsolvedMath ID: 30005404}
\textbf{Source problem code:} catalogue entry
\subsection{Definitions and mathematical statement}
\textbf{Recoverable mathematical content and unresolved target.}
The supplied record restricts to \(n=3\). It begins with the isolated word ``conjecture'' and ends mid-sentence; it supplies no complete conjecture or question. The following is a transcription of its mathematical setup and reported assertions, not a reconstructed open problem.

Write \(\lambda=(\lambda_1,\lambda_2,\lambda_3)\in\Z^3\) for a dominant integral weight of \(\mathfrak{gl}_3\). With the usual upper-triangular convention, dominance means \(\lambda_1\geq\lambda_2\geq\lambda_3\). Let \(W\cong S_3\) be its Weyl group, with simple reflections \(s_1,s_2\) and longest element \(w_0\). The conventional dot action is \(w\cdot\lambda=w(\lambda+\rho)-\rho\), where one may use \(\rho=(2,1,0)\). These conventions explain the notation; the excerpt does not specify the base field or triangular conventions. Let \(M(w\cdot\lambda)\) denote the corresponding Verma module and \(L(\lambda)\) its finite-dimensional simple highest-weight quotient. The excerpt invokes the following terminal portion of the Bernstein--Gelfand--Gelfand resolution:
\[
M(s_1\cdot\lambda)\oplus M(s_2\cdot\lambda)
\longrightarrow M(\lambda)
\longrightarrow L(\lambda)
\longrightarrow 0.
\]
It then reports an exact sequence of sheaves
\[
F_\infty(M(s_1\cdot\lambda))\oplus F_\infty(M(s_2\cdot\lambda))
\longrightarrow F_\infty(M(\lambda))
\longrightarrow F_\infty(L(\lambda))
\longrightarrow 0.
\tag{1}
\]
The functor \(F_\infty\), its source and target categories, its exactness properties, and the maps are not defined in the supplied passage. Thus (1) is retained as a source assertion under omitted hypotheses.

The excerpt says that the sheaves in (1) are Cohen--Macaulay and studies their supports on an undefined space \(E_\infty\), near a chosen form \(f\) and classical forms \(f'\) with the same Galois representation and different refinements. Here a refinement is described as an ordering of the eigenvalues of an operator \(\varphi\); neither that operator nor the hypotheses on the forms and representations are supplied. For a coherent sheaf on a locally Noetherian scheme, the Cohen--Macaulay condition means that every nonzero stalk has depth equal to its dimension as a module. The ambient scheme and coherence hypotheses needed to apply this definition are omitted from the excerpt.

Using a local model attributed to an unidentified reference [3], the record displays the following expression, here reproduced literally in its matrix conditions:
\[
X:=\bigl\{(gB,hB,N)\in
\mathrm{GL}_3/B\times\mathrm{GL}_3/B\times\mathfrak{gl}_3
\;\bigm|\;
 g^{-1}Ng\in\mathfrak b,\quad h^{-1}Ng\in\mathfrak b
\bigr\}.
\]
The notation suggests that \(B\) is a Borel subgroup and \(\mathfrak b\) its Lie algebra, but these choices and the base field are not specified. In particular, the second condition is printed as \(h^{-1}Ng\), not as a conjugation by \(h\); as written, it does not generally define a condition independent of representatives of the flags. It has therefore not been silently corrected.

The record describes \(X\) as a union of pieces \(X_w\), indexed by \(w\in W\), obtained by taking closures of loci where the two flags have relative position \(w\). It calls these pieces ``connected components''; the precise component terminology and relative-position convention require verification. It further writes
\[
Z=G/B\times G/B\subset X_{w_0}
\]
as the locus \(N=0\), without defining \(G\) separately or supplying the conventions needed for this inclusion.

According to the excerpt, \(X\) models the image of \(E_\infty\) in an undefined space \(X_\infty\) at \(f\), and the particular choice of \(f\) corresponds to
\[
x_0=(B,B,0)\in X_{w_0}.
\]
It asserts that \(X_{w_0}\) is not Gorenstein at \(x_0\), that the other components are smooth, and that \(Z\) is smooth. The choice of \(f\), the local-model comparison, and all these geometric assertions depend on omitted context and are not independently verified here.

Finally, the passage associates \(F_\infty(M(\lambda))\) with \(\mathcal O_{X_{w_0},x_0}\), and associates the other three sheaves in (1) with the local rings at \(x_0\) of
\[
X_{s_2s_1},\qquad X_{s_1s_2},\qquad Z.
\]
It phrases these associations as identifications of supports with local rings. The precise distinction between a support and a module over its local ring, and the individual correspondence for the remaining three sheaves, are not explicitly supplied. The fragment then ends with ``Those last three being smooth, the corresponding sheaves are'', leaving its conclusion missing.

\textbf{Statement limitation.}
No complete mathematical problem can be stated from this record without adding a missing target, hypotheses, and corrections to the local-model formula. All assertions above remain attributed fragmentary exposition requiring human mathematical review.
\subsection{Short English statement}
For \(n=3\), the supplied fragment discusses a BGG sequence, its images under an undefined functor \(F_\infty\), and proposed descriptions of sheaf supports using a Steinberg local model. It ends before its conclusion and contains unresolved notation and hypotheses. The supplied registry labels the item open, but also marks its statement unrecoverable and describes the passage as proof exposition rather than an open problem; no complete conjecture can be recovered here.
\subsection{Sources}
\begin{itemize}
\item Supplied local UnsolvedMath registry record, permanent ID \(30005404\), exact source problem code \texttt{catalogue entry}, title ``Local-Model Description of Shimura Varieties'', in the Oberwolfach Reports Open Problems collection. The proposed year is 2023; no proposer is named.
\item Supplied report citation: \emph{Arithmetic of Shimura Varieties} (2023), Oberwolfach Reports, DOI 10.4171/owr/2023/5, pp. 261--326. \url{https://doi.org/10.4171/owr/2023/5}. The record attributes the material to workshop reports and their open-problem discussions; the supplied passage itself is incomplete.
\item Supplied literature source, labeled ``original OWR report'': \url{https://oa.tib.eu/renate/server/api/core/bitstreams/ae468b44-47d2-4568-9d88-ffcb4023f8d4/content}.
\item Supplied literature source, labeled ``related patching results'': \url{https://arxiv.org/abs/2406.01129}. No authors, title, or precise supporting result are supplied for this link.
\item Literature assessment dated exactly \texttt{2026-08-22}: supplied dated curation, not independently verified. It labels the item open but explains that the record ends mid-argument after discussing BGG images and Steinberg-variety supports, leaving no complete claim whose truth can be certified. The supplied background requests independent verification of time-sensitive status and cited claims.
\item The supplied statement-verification field describes the official-report passage as proof and local-model exposition rather than an open mathematical problem, and the statement-status field is \texttt{unrecoverable}. These are supplied assessments, not findings from an independent source visit. The internal reference [3] is not identified by the supplied record.
\end{itemize}
\end{document}
